\chapter{Introduction}
\label{ch:introduction}

\section{Background}
\label{sec:background}

Road traffic accidents remain one of the leading causes of preventable
death and injury in the modern world. According to the World Health
Organization's \emph{Global Status Report on Road Safety}, approximately
1.19 million people lose their lives in road crashes every year, and
between twenty and fifty million more suffer non-fatal injuries that
often result in long-term disability. Driver-related factors account
for the overwhelming majority of these incidents, with human error
contributing to roughly ninety percent of all recorded crashes.
Historically, attention has been concentrated on the most visible
contributors such as speeding, alcohol intoxication, and fatigue.
However, a growing body of epidemiological and behavioral research
now implicates \emph{acute psychological stress} as a silent and
under-recognized risk factor behind a substantial share of these
events \citep{sharma2012driver}.

Driving is a cognitively demanding task. It requires sustained visual
attention, quick decision making, and fine motor control all at once.
When a driver is under acute psychological stress, a number of things
change: heart rate variability drops, breathing becomes shallower,
muscle tension builds in the shoulders and neck, and the sympathetic
branch of the autonomic nervous system starts dominating cardiovascular
regulation. These changes narrow the visual field (``tunnel vision''),
slow reaction times, reduce inhibitory control, and push drivers toward
impulsive maneuvers --- tailgating, abrupt lane changes, failure to
yield. Epidemiological studies have consistently found that high
self-reported driving stress is linked to a two- to threefold increase
in crash risk, especially during rush-hour commutes and long-haul
driving.

The sources of driver stress are diverse. External factors include
traffic congestion, bad weather, and complex road geometry. Internal
factors include time pressure, interpersonal conflict, workplace
deadlines, and pre-existing anxiety. Commercial drivers (taxi,
truck, and ride-share) are especially vulnerable because long hours
behind the wheel combine with economic pressure and irregular
schedules. Younger and novice drivers are another at-risk group
because they tend to have higher baseline anxiety while also lacking
the automated driving skills that make experienced driving less
mentally taxing.

Despite the clear public-health significance of this problem,
contemporary vehicles offer almost no countermeasures aimed at
driver stress. Modern driver-assistance systems address perceptual
lapses (lane keeping, forward collision warning, adaptive cruise
control) and gross behavioral impairment (drowsiness detection,
alcohol interlocks), but they're essentially blind to the driver's
emotional and autonomic state. That gap is what this report addresses:
we close it using wearable sensors that many drivers already carry on
their wrists and machine-learning models that run entirely on-device
within the power and memory budget of a smartwatch.

\section{Motivation}
\label{sec:motivation}

Real-time, software-based driver stress detection is feasible today
for reasons that didn't exist a decade ago.

\textbf{Consumer wearables now carry sensor-grade hardware.}
Devices like the Samsung Galaxy Watch, Google Pixel Watch, Fitbit
Sense, and Apple Watch routinely expose photoplethysmography (PPG),
skin temperature, three-axis acceleration, and on selected models
electrodermal activity (EDA). That sensor complement was, until
recently, confined to laboratory-grade equipment like the Empatica
E4 or the RespiBAN chest band. Now it ships on devices worn by
hundreds of millions of people for twelve-plus hours a day. The
WESAD dataset \citep{schmidt2018wesad}, which underpins this project,
was collected using exactly this wrist-wearable configuration, so
its findings carry over to mass-market devices without modification.

\textbf{On-device machine learning has matured.}
TensorFlow Lite \citep{tflite2023}, ONNX Runtime Mobile, and Core ML
give developers production-ready toolchains for running quantized
neural networks on battery-powered hardware. Modern smartwatch
system-on-chips have hardware-accelerated inference paths that execute
small recurrent models in milliseconds while drawing only a few
milliwatts. In practice, this means stress detection can run
\emph{entirely on the wrist}: no cellular link, no cloud round-trip,
no continuous upload of biometric data. That matters a lot for
automotive use cases where network coverage is patchy and privacy
expectations are high.

\textbf{Low-cost BLE microcontrollers make the actuation side tractable.}
Boards like the Arduino Nano 33 BLE are inexpensive, widely available,
and capable of driving relays, LEDs, fans, and infotainment inputs while
speaking a standard GATT profile to any nearby wearable. An end-to-end
prototype -- wearable sensing to cabin actuation -- no longer requires
bespoke hardware or any invasive modification to the vehicle.

The \emph{graduated} nature of the intervention is worth highlighting
separately. Alcohol interlocks make a binary call: ignition allowed or
blocked. A stress-aware cabin system can do more than that. A mild
stress reading can trigger a subtle lighting change and a slight audio
volume reduction. A moderate reading can additionally switch the LED
color to a calming blue and add soft background music. Only a critical
reading triggers an overt visual and auditory alert. This kind of
graduated, non-punitive response is more likely to be accepted by
drivers than a hard override, and it doesn't compete with existing
driver-assistance systems.

\section{Project Scope}
\label{sec:scope}

This project is deliberately \textbf{software-focused}. No new hardware
has been designed. The sensing hardware (wrist PPG, EDA, skin temperature,
and accelerometer) is available in off-the-shelf consumer wearables.
The actuation hardware (LEDs, PWM-driven cooling fan, and piezo speaker)
is directly supported by an unmodified Arduino Nano 33 BLE development
board. The contribution of the project lies
entirely in (i) the machine-learning model and its training pipeline,
(ii) the Wear OS Kotlin application that runs the model in real time,
(iii) the BLE protocol layer that conveys stress estimates to the
vehicle, and (iv) the Arduino firmware that translates those estimates
into cabin actions.

The scope of the scientific evaluation is equally bounded. All training
and testing is performed on the publicly available WESAD dataset. No
naturalistic driving data has been collected; we do not claim that the
system has been validated \emph{in vivo} on real drivers, and this
limitation is revisited explicitly in Chapter~\ref{ch:conclusion}.
The system is, however, architected in a way that does not depend on
laboratory assumptions -- it expects streaming sensor data, tolerates
missing samples, and emits its estimates on a fixed thirty-second
cadence suitable for deployment outside the laboratory.

The scope is also bounded on the safety side. The adaptive cabin
control system is \emph{advisory}: it changes lighting, fan speed,
and audio profile, but it does not touch steering, braking, or any
other primary vehicle control. That keeps the prototype firmly in the
regulatory space of an aftermarket comfort accessory, not an automated
driving component, which is consistent with the graduated-intervention
approach described above.

\section{Problem Statement}
\label{sec:problem}

Despite extensive prior work on both stress physiology and wearable
biosensing, four concrete gaps remain unaddressed in the literature
and the market:

\begin{enumerate}
  \item \textbf{No real-time multimodal wearable stress classifier
  targeted at drivers.} The most successful stress classifiers
  published on WESAD use offline, PC-class models that are neither
  quantized nor latency-budgeted for a smartwatch, and many of them
  rely on chest-band signals (respiration, ECG) that wrist wearables
  cannot acquire.

  \item \textbf{No adaptive cabin response system driven by
  physiological state.} Commercial vehicles offer static ``comfort
  profiles'' selected manually by the driver; none of them respond
  automatically to the driver's real-time autonomic state, and none
  of them gradate the response across distinct stress levels.

  \item \textbf{Existing solutions are not wearable-constrained.}
  State-of-the-art stress detection pipelines frequently assume access
  to full-resolution raw signals at 64--700~Hz and to tens of megabytes
  of RAM. Consumer smartwatches offer neither. Bridging that gap
  requires aggressive signal downsampling, a small model footprint,
  and quantisation-aware conversion tooling.

  \item \textbf{No integrated end-to-end open system.} To our
  knowledge, no published work combines wearable sensing, on-device
  inference, a documented BLE protocol, and a vehicle-side actuator
  controller in a single openly described prototype.
\end{enumerate}

These gaps together define the problem addressed in this project.

\section{Objectives}
\label{sec:objectives}

To close the gaps identified above, this project pursues the
following five concrete objectives:

\begin{enumerate}
  \item \textbf{Design and train a four-class driver stress classifier}
  (Calm / Mild / Moderate / Critical) based on a bidirectional LSTM
  \citep{graves2005bilstm} with an additive attention layer
  \citep{vaswani2017attention}, using only wrist-accessible signals
  from the WESAD dataset \citep{schmidt2018wesad}.

  \item \textbf{Achieve accuracy above 90\% and macro-F1 above 0.90}
  on a held-out test split while keeping the model size compatible
  with on-device inference (target $\le 80$~KB after TensorFlow Lite
  conversion).

  \item \textbf{Implement a Wear OS Kotlin application} that performs
  real-time sensor collection, windowing, normalization, and on-device
  TensorFlow Lite inference \citep{tflite2023} with an end-to-end
  latency below fifty milliseconds.

  \item \textbf{Specify and implement a compact BLE protocol}
  (a twelve-byte GATT characteristic) carrying timestamp, stress level,
  and confidence, and connect it to an Arduino Nano 33 BLE that
  implements the vehicle-side finite state machine.

  \item \textbf{Demonstrate a graduated adaptive cabin response}
  consisting of four differentiated cabin profiles (Calm, Mild,
  Moderate, Critical), a waiting state for connection loss, and a
  sixty-second BLE timeout enforced at the Arduino.
\end{enumerate}

Each subsequent chapter of this report addresses one or more of
these objectives directly, and Chapter~\ref{ch:conclusion} revisits
them in order to demonstrate that each has been met by the
implemented system.

\section{Connection to Patent ACN1408}
\label{sec:acn1408}

This work extends a prior invention by the same author, filed as an
Indian Provisional Patent Application under the number ACN1408
\citep{acn1408}. The earlier invention -- commercially referred to
as \emph{AlcoWatch} -- addresses the problem of \textbf{alcohol-related
driving accidents} by using a wrist wearable to estimate blood alcohol
concentration (BAC) from physiological correlates (PPG, EDA, skin
temperature) and by using a relay-controlled Arduino module to block
vehicle ignition when BAC exceeds the legal limit of 0.08~g/dL. In
that system the actuation is binary (ignition ALLOW / ignition BLOCK),
the failure mode is conservative (ignition is blocked by default),
and the model is regressive (it outputs a scalar BAC estimate).

The present project extends ACN1408 along two axes. On the
intervention side, it moves from a \emph{binary safety interlock}
(ignition block) to a \emph{graduated comfort response} (cabin
profile). On the ML side, it moves from \emph{BAC regression} to
\emph{stress classification} with four ordinal classes. The two
systems are architecturally compatible: they use the same class of
wearable hardware, the same class of Arduino microcontroller, and a
very similar BLE service structure. Chapter~\ref{ch:architecture} provides a direct side-by-side
comparison, and Chapter~\ref{ch:conclusion} sketches how the two could
be unified into a single multi-purpose wearable safety platform
in future work.
